
В данной работе рассматривается задача проверки того, что наблюдаемое
исполнение распределённой программы согласуется с её высокоуровневой формальной
спецификацией на языке TLA+. В отличие от \emph{уточнения} (refinement), где
требуется доказать, что \emph{все} поведения реализации допустимы
спецификацией, решается более практическая задача: проверить, что
\emph{конкретная} записанная трасса исполнения может быть интерпретирована как
некоторое поведение спецификации (\emph{trace validation}, \emph{model-based
trace-checking})

\subsubsection{Математическая модель спецификации.}

\begin{definition}\label{spec}
Пусть $Vars=\langle v_1,\dots,v_k\rangle$ --- конечный кортеж переменных.
Обозначим множество имён переменных через
\[
\mathit{VNames} \triangleq \{v_1,\dots,v_k\}.
\]
Пусть $\mathcal{D}$ --- множество допустимых значений (домен значений), которые
могут принимать переменные спецификации (строки, множества, функции, записи и
т.д. в смысле математических объектов).

\emph{Состоянием} спецификации называется тотальная функция
\[
s:\mathit{VNames}\to \mathcal{D},
\]
которая каждой переменной $v\in\mathit{VNames}$ сопоставляет её значение $s(v)$.

Пусть $\mathcal{S}\triangleq \mathcal{D}^{\mathit{VNames}}$ --- множество всех состояний.
\emph{Предикатом на состояниях} называется подмножество $I\subseteq\mathcal{S}$
(или эквивалентно характеристическая функция $I:\mathcal{S}\to\{\mathsf{true},\mathsf{false}\}$).
\emph{Отношением переходов} называется подмножество
\[
N \subseteq \mathcal{S}\times \mathcal{S},
\]
элемент $(s,s')\in N$ интерпретируется как допустимый шаг из состояния $s$ в
состояние $s'$.

\emph{Спецификацией} (в форме машины состояний) называется пара
\[
\mathit{Spec}\triangleq (I,N),
\]
где $I\subseteq\mathcal{S}$ задаёт множество начальных состояний, а
$N\subseteq\mathcal{S}\times\mathcal{S}$ задаёт отношение ``следующее состояние''.

\emph{Конечным поведением} длины $n\ge 0$ называется последовательность состояний
\[
\sigma=\langle s_0,s_1,\dots,s_n\rangle,\qquad s_i\in\mathcal{S}.
\]
Говорят, что $\sigma$ \emph{удовлетворяет} спецификации $\mathit{Spec}=(I,N)$,
если
\[
s_0\in I\ \wedge\ \forall i\in\{0,\dots,n-1\}:\ (s_i,s_{i+1})\in N.
\]
Множество всех конечных поведений, удовлетворяющих спецификации, обозначим
\[
\mathcal{B}_f(\mathit{Spec}) \triangleq
\bigcup_{n\ge 0}\left\{
\langle s_0,\dots,s_n\rangle \in \mathcal{S}^{n+1}\ \middle|\
s_0\in I \ \wedge\ \forall i<n:\ (s_i,s_{i+1})\in N
\right\}.
\]
\end{definition}

Пусть задана TLA+ спецификация $\mathit{Spec}$ над конечным набором переменных
$Vars = \langle v_1, \dots, v_k \rangle$ (например, $Vars = \langle rmState,
tmState, tmPrepared, msgs$).


Спецификация задаёт
множество допустимых (конечных) поведений как последовательностей состояний:
\[
  S \triangleq \{ \sigma = \langle s_0, s_1, \dots, s_n \rangle \mid \sigma \models \mathit{Spec} \}.
\]
Здесь состояние $s_i$ есть отображение $Vars \to \mathcal{V}$, где
$\mathcal{V}$ --- домен значений TLA+.

Для дальнейшего удобно фиксировать, что $\mathit{Spec}$ имеет вид
\[
\mathit{Spec} \triangleq \mathit{Init} \wedge \Box[\mathit{Next}]_{Vars},
\]
где $\mathit{Init}$ задаёт начальные состояния, а $\mathit{Next}$ --- отношение
перехода между соседними состояниями \cite{LamportSpecifyingSystems2002}.

\paragraph{Модель наблюдаемой трассы.}
Пусть программа при исполнении порождает лог $\tau$ длины $m$,
состоящий из записей $\tau[1],\dots,\tau[m]$. Каждая запись фиксирует:
(i) имя события (\texttt{event}) и параметры (\texttt{event\_args});
(ii) изменения некоторых переменных спецификации в виде \emph{обновлений}
(update operations), а не полных значений.

Пусть $Obs \subseteq \{v_1,\dots,v_k\}$ --- множество переменных, которые
мы действительно наблюдаем (логируем). Трасса обычно \emph{неполная}:
часть переменных не записывается вовсе, а для записываемых переменных
фиксируются лишь приращения (\emph{updates}). Это принципиально важно:
неполная трасса не задаёт единственную последовательность состояний,
а задаёт \emph{множество} возможных последовательностей.

Формально определим отношение совместимости состояния $s$ и записи лога $\tau[i]$.
Пусть $\mathsf{Apply}(\tau[i], s)$ --- множество состояний $s'$, которые могут
получиться из $s$ применением всех обновлений, присутствующих в $\tau[i]$,
к соответствующим переменным из $Obs$ (обновления задаются операторами
типа \texttt{Update}, \texttt{AddElement} и т.п.).
Если некоторой переменной $v \in Obs$ нет в записи $\tau[i]$,
то её значение может \emph{либо} не изменяться (частый случай),
\emph{либо} оставаться неограниченным (если мы допускаем, что трасса
неполна даже по отношению к изменившимся переменным). В данной работе
используется практичная конвенция: \emph{если обновления для переменной не
указаны, то переменная не изменяется}. Этот выбор соответствует
«логированию точек линеаризации» (linearization points), когда лог пишется
после завершения атомарного шага \cite{CirsteaKuppeLoillierMerz2024}.

Тогда определим множество \emph{поведений, совместимых с трассой}:
\[
T \triangleq \left\{ \sigma=\langle s_0,\dots,s_m\rangle \ \middle|\
\forall i \in 1..m: s_i \in \mathsf{Apply}(\tau[i], s_{i-1}) \right\}.
\]
Заметим, что $T$ может быть большим из-за не наблюдаемых переменных и/или
неполной информации.

\paragraph{Задача проверки трассы.}
Мы говорим, что трасса $\tau$ \emph{совместима} со спецификацией $\mathit{Spec}$,
если существует хотя бы одно поведение, которое одновременно:
(1) удовлетворяет спецификации; (2) согласуется с наблюдаемой трассой.
То есть требуется проверить:
\[
S \cap T \neq \varnothing.
\]
Именно это условие является целевым критерием валидации трасс
\cite{HowardTraceChecking2011,CirsteaKuppeLoillierMerz2024}.

Важно подчеркнуть, что мы \emph{не} проверяем включение $T \subseteq S$,
то есть не доказываем уточнение. Для неполных логов это условие слишком сильное:
выбор значений не наблюдаемых переменных мог бы породить поведения,
которые программа никогда не реализует, но которые формально «разрешены»
логом как множеством $T$.

\paragraph{Проблема атомарности и стуттеринга.}
Ещё одна тонкость заключается в несовпадении «крупности» шагов реализации и
шагов спецификации. Один атомарный переход спецификации может соответствовать:
\begin{itemize}
  \item одному «логическому шагу» программы (идеальный случай);
  \item нескольким низкоуровневым операциям, которые в логе объединяются в один шаг;
  \item либо, наоборот, одному шагу программы могут соответствовать \emph{несколько}
        переходов спецификации (или стуттеринг-шага спецификации).
\end{itemize}
В TLA+ эта проблема естественным образом трактуется через
\emph{stuttering equivalence}: допускаются шаги, которые не меняют значения
переменных $Vars$ (или выбранного поднабора), но «проталкивают» исполнение
вперёд \cite{LamportSpecifyingSystems2002}. На уровне трасс-валидации это приводит к тому,
что запись $\tau[i]$ может соответствовать нулю или нескольким переходам $\mathit{Next}$,
а значит, формулу совместимости нужно корректно редуцировать к модели TLC.

\paragraph{Редукция к constrained model checking.}
Покажем, что проверка $S \cap T \neq \varnothing$ сводится к проверке
существования поведения у \emph{ограниченной спецификации}
(\emph{trace specification}), построенной из $\mathit{Spec}$ и трассы $\tau$.

Введём дополнительную переменную $l$ (line index), которая указывает на номер
текущей записи трассы, и определим:
\[
\mathit{TraceInit} \triangleq (l = 1),
\]
а также отношение перехода $\mathit{TraceStep}$, которое:
(1) применяет обновления из $\tau[l]$ к наблюдаемым переменным;
(2) проверяет, что выполнен некоторый допустимый переход спецификации;
(3) увеличивает $l$ на 1.
Интуитивно:
\[
\mathit{TraceStep} \triangleq
\mathsf{UpdateFromTrace}(\tau[l]) \ \wedge\ \mathit{Next} \ \wedge\ (l' = l + 1).
\]
Тогда определим \emph{ограниченную спецификацию}:
\[
\mathit{Spec}_\tau \triangleq (\mathit{Init} \wedge \mathit{TraceInit})
\ \wedge\ \Box[\mathit{TraceStep}]_{\langle Vars, l\rangle}.
\]

\begin{theorem}[Корректность редукции]
Трасса $\tau$ совместима со спецификацией $\mathit{Spec}$
тогда и только тогда, когда у спецификации $\mathit{Spec}_\tau$
существует поведение длины $m$ (то есть достигается $l = m+1$).
\end{theorem}

\begin{proof}
($\Rightarrow$) Пусть $S \cap T \neq \varnothing$.
Тогда существует $\sigma=\langle s_0,\dots,s_m\rangle$, такое что
$\sigma \models \mathit{Spec}$ и $\sigma \in T$.
Рассмотрим расширение этой последовательности состояниями переменной $l$:
определим $l_0 = 1$, $l_i = i+1$.
Поскольку $\sigma \models \mathit{Spec}$, каждое $\langle s_{i-1}, s_i\rangle$
удовлетворяет $\mathit{Next}$ (с учётом $[\cdot]_{Vars}$).
Поскольку $\sigma \in T$, переход $s_{i-1}\to s_i$ согласуется с обновлениями
из $\tau[i]$, то есть выполняется $\mathsf{UpdateFromTrace}(\tau[i])$.
Следовательно, каждый расширенный шаг удовлетворяет $\mathit{TraceStep}$,
а начальное состояние удовлетворяет $\mathit{Init}\wedge(l=1)$.
Итак, получено поведение спецификации $\mathit{Spec}_\tau$.

($\Leftarrow$) Пусть существует поведение $\langle (s_0,l_0),\dots,(s_m,l_m)\rangle$
спецификации $\mathit{Spec}_\tau$, где $l_0=1$ и $l_m=m+1$.
Тогда по определению $\mathit{Spec}_\tau$ проекция $\langle s_0,\dots,s_m\rangle$
удовлетворяет $\mathit{Spec}$, поскольку $\mathit{TraceStep}$ включает $\mathit{Next}$,
а $\mathit{Init}$ сохраняется в начальном состоянии.
Кроме того, поскольку на каждом шаге выполняется $\mathsf{UpdateFromTrace}(\tau[l])$,
получаем, что $\langle s_0,\dots,s_m\rangle \in T$.
Значит, $\langle s_0,\dots,s_m\rangle \in S \cap T$, т.е. $S \cap T \neq \varnothing$.
\end{proof}

Таким образом, задача валидации трассы редуцируется к задаче
\emph{поиска хотя бы одного поведения} у $\mathit{Spec}_\tau$.
Эта постановка естественным образом решается методом модельной проверки,
где трасса выступает ограничением на пространство состояний
(\emph{constrained model checking}) \cite{CirsteaKuppeLoillierMerz2024,HowardTraceChecking2011}.

\paragraph{Замечание о (не)полноте.}
Если трасса неполна, то $\mathit{Spec}_\tau$ может допускать поведения,
которые соответствуют трассе, но не реализуемы реальной программой.
Следовательно, успех валидации даёт \emph{доказательство совместимости},
но не гарантирует отсутствия ошибок. Напротив, провал валидации означает,
что наблюдаемое исполнение не может быть объяснено спецификацией при
принятых соглашениях о логировании и атомарности, что является сильным
сигналом о наличии ошибки в реализации или в инструментировании
\cite{CirsteaKuppeLoillierMerz2024}.

\subsection{Проблема верификации соответствия трасс 2}

Эта работа решает задачу \emph{валидации трасс} (trace validation) исполнения
распределённой программы относительно высокоуровневой спецификации на TLA+.
Интуитивно мы хотим ответить на вопрос:

\begin{quote}
``Может ли наблюдаемая трасса исполнения быть объяснена некоторым корректным
поведением спецификации?''
\end{quote}

В отличие от \emph{доказательства уточнения} (refinement), где требуется показать,
что \emph{все} возможные поведения реализации допустимы спецификацией,
здесь проверяется существование хотя бы одного допустимого поведения,
совместимого с наблюдаемой трассой
(\emph{model-based trace-checking}) \cite{HowardTraceChecking2011,CirsteaKuppeLoillierMerz2024}.

\paragraph{Нотация и соглашения.}
В тексте используются стандартные математические обозначения.
\begin{itemize}
    \item Символ $\triangleq$ означает \emph{определение по соглашению}
    (``равно по определению''). То есть запись $A \triangleq B$ читается как
    ``определим $A$ как $B$''.
    \item $\langle x_0, x_1, \dots, x_n \rangle$ --- последовательность (sequence).
    \item Для множества $X$ запись $x \in X$ означает ``$x$ принадлежит $X$''.
    \item Для отображения (функции) $s : Vars \to \mathcal{V}$
    запись $s(v)$ означает значение переменной $v$ в состоянии $s$.
    \item Для множества последовательностей $A$ и $B$ их пересечение
    $A \cap B$ — множество последовательностей, принадлежащих обоим.
\end{itemize}

\paragraph{1. Спецификация как множество поведений.}
Пусть дана TLA+ спецификация $\mathit{Spec}$, в которой объявлены переменные
состояния:
\[
Vars \triangleq \langle v_1, v_2, \dots, v_k \rangle.
\]
Тогда \emph{состояние} $s$ — это функция, сопоставляющая каждой переменной
значение из домена значений:
\[
s : Vars \to \mathcal{V}.
\]
Например, в спецификации TwoPhaseCommit:
\[
Vars = \langle rmState, tmState, tmPrepared, msgs \rangle.
\]

\emph{Поведение} (behavior) — конечная последовательность состояний:
\[
\sigma \triangleq \langle s_0, s_1, \dots, s_n \rangle.
\]

Спецификация $\mathit{Spec}$ задаёт множество всех допустимых поведений:
\[
S \triangleq \{ \sigma \mid \sigma \models \mathit{Spec} \}.
\]
Здесь $\sigma \models \mathit{Spec}$ означает: ``поведение $\sigma$
удовлетворяет формуле спецификации'' (как в стандартной семантике TLA+).

Часто спецификацию записывают в виде:
\[
\mathit{Spec} \triangleq \mathit{Init} \ \wedge\ \Box[\mathit{Next}]_{Vars},
\]
где:
\begin{itemize}
    \item $\mathit{Init}$ — предикат на начальное состояние $s_0$;
    \item $\mathit{Next}$ — предикат на пару соседних состояний $(s_i, s_{i+1})$,
    то есть отношение перехода;
    \item $\Box[\mathit{Next}]_{Vars}$ допускает \emph{стуттеринг-шаги}:
    либо выполняется $\mathit{Next}$, либо состояние не меняется по переменным $Vars$.
\end{itemize}
Такое представление является каноническим для TLA+ \cite{LamportSpecifyingSystems2002}.

\paragraph{2. Что такое ``трасса'' в нашей постановке.}
Пусть программа при исполнении порождает лог (trace) длины $m$:
\[
\tau \triangleq \langle \tau[1], \tau[2], \dots, \tau[m] \rangle,
\]
где каждая $\tau[i]$ — запись (например, строка NDJSON), содержащая:
\begin{itemize}
    \item имя события (например, \texttt{RMPrepare}, \texttt{TMCommit});
    \item список обновлений некоторых переменных спецификации
          (например, \texttt{rmState}, \texttt{tmPrepared}).
\end{itemize}

Важно: трасса обычно \emph{не хранит полные значения переменных},
а хранит \emph{обновления} (updates). Поэтому одна и та же трасса может быть
совместима с несколькими разными последовательностями состояний.

\paragraph{3. Множество поведений, допускаемых трассой (из-за неполноты).}
Чтобы формализовать неполноту, введём множество логируемых переменных:
\[
Obs \subseteq \{v_1,\dots,v_k\}.
\]
Например, $Obs = \{rmState, tmState, tmPrepared, msgs\}$,
или (в более экономном режиме) $Obs = \{tmState, tmPrepared\}$.

Далее определим, что означает ``запись $\tau[i]$ совместима с переходом
из состояния $s$ в состояние $s'$''.

\paragraph{Отношение обновления по записи трассы.}
Определим бинарное отношение:
\[
(s, s') \in \mathsf{UpdateByLine}(\tau[i])
\]
как ``состояние $s'$ получается из $s$ применением обновлений,
заданных в $\tau[i]$, к переменным из $Obs$''.

Смысл:
\begin{itemize}
    \item если в $\tau[i]$ для переменной $v \in Obs$ указано обновление,
    то $s'(v)$ обязано быть равно результату применения этого обновления к $s(v)$;
    \item если переменная $v \in Obs$ \emph{не упомянута} в $\tau[i]$,
    мы принимаем практичное соглашение:
    \[
      s'(v) = s(v)
    \]
    (то есть ``не залогировано'' интерпретируется как ``не менялось'').
\end{itemize}

Это соглашение соответствует идее линейризационных точек:
мы логируем момент после завершения атомарного шага, поэтому все изменения,
важные для спецификации, должны появляться в логе \cite{CirsteaKuppeLoillierMerz2024}.

\paragraph{Множество поведений, совместимых с трассой.}
Используя $\mathsf{UpdateByLine}$, определим множество последовательностей состояний,
которые \emph{можно построить по трассе}:
\[
T \triangleq \left\{
\langle s_0, s_1, \dots, s_m \rangle \ \middle|\
\forall i \in 1..m:\ (s_{i-1}, s_i) \in \mathsf{UpdateByLine}(\tau[i])
\right\}.
\]
Это множество $T$ является ``семантикой трассы''.
Если трасса неполна (не все переменные фиксируются, не все значения известны),
то $T$ обычно содержит много разных последовательностей состояний.

\paragraph{4. Определение совместимости трассы и спецификации.}
Теперь можно строго определить основную задачу.

\begin{definition}[Совместимость трассы со спецификацией]
Трасса $\tau$ \emph{совместима} со спецификацией $\mathit{Spec}$, если
\[
S \cap T \neq \varnothing.
\]
\end{definition}

То есть существует хотя бы одно поведение $\sigma$ такое, что:
\begin{itemize}
    \item $\sigma$ удовлетворяет спецификации ($\sigma \in S$),
    \item и одновременно $\sigma$ согласуется с трассой ($\sigma \in T$).
\end{itemize}

Если пересечение непусто, мы говорим, что трасса \textbf{принята} (accepted).
Если пересечение пусто, трасса \textbf{отклонена} (rejected): наблюдаемое исполнение
не может быть объяснено спецификацией при заданных правилах интерпретации лога.

\paragraph{5. Почему мы не проверяем refinement.}
Можно было бы попытаться проверить условие $T \subseteq S$,
то есть ``все поведения, допускаемые трассой, разрешены спецификацией''.
Но это условие почти всегда слишком сильное, потому что:
\begin{itemize}
    \item трасса неполна и допускает ``лишние'' варианты значений для не наблюдаемых
    переменных;
    \item трасса — это выборка, она не описывает все возможные исполнения программы.
\end{itemize}
Поэтому критерий $S \cap T \neq \varnothing$ является реалистичной и широко
используемой постановкой model-based trace-checking \cite{HowardTraceChecking2011}.

\paragraph{6. Дополнение: связь событий лога и действий спецификации.}
В реальной системе запись $\tau[i]$ обычно содержит поле \texttt{event},
которое связывается с действием $\mathit{Next}$ в TLA+.
Например, \texttt{event = "TMRcvPrepared"} соответствует действию
$TMRcvPrepared(r)$, где параметр $r$ берётся из \texttt{event\_args}
или выбирается недетерминированно, если он не указан.
Это связывание формирует дополнительные ограничения поверх $\mathsf{UpdateByLine}$,
сужая множество $T$ до ``тех'' переходов, которые соответствуют заявленным событиям.
Именно это затем реализуется в Trace Specification для TLC
\cite{CirsteaKuppeLoillierMerz2024}.

